\chapter{Focus and keyboard navigation}
\label{ch:focus}

\begin{aside}
A UI with more than one focusable region needs a way to say
``input goes \emph{here} now''. \maya{} models focus as state in
your \code{Model}, not as a global property of the runtime.
That sounds restrictive; it is liberating. Focus becomes data,
flows through \code{update} like any other state, and is
testable without a runtime.
\end{aside}

\section{Why focus is not a global}

Most GUI frameworks store focus as a property of the widget
tree: a pointer to the ``focused widget'', set by mouse clicks
and Tab presses, read by event dispatchers. The pattern works
but it has costs:

\begin{itemize}[leftmargin=*]
  \item \textbf{Hidden state.} Focus is not visible in your
    model; it is owned by the framework. Diagnosing ``why did
    keystrokes go to the wrong place?'' means reading framework
    internals.
  \item \textbf{Time-travel breaks.} Restoring an old model does
    not restore the old focus; the framework's pointer survives.
  \item \textbf{Concurrency hazards.} Mouse and keyboard handlers
    might race on the focus pointer.
\end{itemize}

The Elm fix: focus is just a field in your \code{Model}. The
runtime never touches it; you decide who has focus by setting
the field in \code{update}.

\subsection*{The simplest form}

\begin{lstlisting}
struct Model {
    int focus;    // 0 = left, 1 = right
    LeftPanel  left;
    RightPanel right;
};
\end{lstlisting}

The \code{view} renders both panels but highlights only the
focused one. The keyboard handler routes input based on
\code{focus}. \code{Tab} flips \code{focus}.

That is the whole pattern. There is no \code{focus\_widget()}
call.

\section{The FocusManager helper}

For larger apps, naming focus targets by integers gets old.
\file{maya/include/maya/core/focus.hpp} ships a small helper:

\begin{lstlisting}
class FocusManager {
public:
    using Id = uint32_t;

    Id  register_target();
    void focus(Id);
    void focus_next();
    void focus_prev();
    bool is_focused(Id) const;
    Id   current() const;

private:
    std::vector<Id> order_;
    Id current_ = 0;
};
\end{lstlisting}

Each focusable region calls \code{register\_target()} at
construction to get a stable ID. \code{focus\_next} cycles
forward, \code{focus\_prev} cycles back.

The manager itself is plain data --- a vector and an integer.
You can store one in your \code{Model} and read it from
\code{view}.

\subsection*{Putting focus in the view}

\begin{lstlisting}
Element render_panel(const Panel& p, bool focused) {
    return box(
        text(p.label) | (focused ? Bold : Dim),
        text(p.body))
        | border<focused ? Heavy : Single>;
}
\end{lstlisting}

The same panel renders differently when focused. The decision is
the caller's; no global ``am I focused?'' query.

\subsection*{Keyboard routing}

\begin{lstlisting}
static std::pair<Model, Cmd<Msg>>
update(Model m, Msg msg) {
    return std::visit(overload{
        [&](KeyPress k) {
            if (k.key == Tab) {
                m.focus.focus_next();
                return std::pair{m, Cmd<Msg>{}};
            }
            // Route to the focused panel.
            if (m.focus.is_focused(m.left.id))
                return route_to_left(m, k);
            return route_to_right(m, k);
        },
        // ...
    }, msg);
}
\end{lstlisting}

The handler decides who gets the key based on
\code{focus.current()}. The focused panel's own \code{update}
processes it.

\section{Tab order}

\code{FocusManager} maintains the registration order. By default
\code{focus\_next} cycles through it. You can override:

\begin{lstlisting}
mgr.set_order({sidebar_id, main_id, search_id});
\end{lstlisting}

Or skip-disabled-elements:

\begin{lstlisting}
void focus_next_enabled(FocusManager& m, const Model& model) {
    do {
        m.focus_next();
    } while (!model.is_target_enabled(m.current())
             && m.current() != initial_id);
}
\end{lstlisting}

The pattern: focus order is a function of your model, not the
manager. The manager is a bookkeeping helper; the rules are
yours.

\section{Mouse focus}

A click can change focus. The hit test
(Chapter~\ref{ch:elements}) returns the deepest element under
the cursor; if that element has an associated focus ID, set it:

\begin{lstlisting}
[&](MouseEvent me) {
    if (me.kind == MouseKind::Press && me.button == 0) {
        if (auto id = hit_to_focus_id(m, me.x, me.y)) {
            m.focus.focus(*id);
        }
    }
    return std::pair{m, Cmd<Msg>{}};
}
\end{lstlisting}

The mapping from rect to focus ID lives in your model. The
runtime does not know which boxes are focusable; you do.

\section{Modal focus}

A modal (a dialog that captures all input) is just a focus stack:

\begin{lstlisting}
struct Model {
    FocusManager focus;
    std::vector<FocusManager> modal_stack;
};
\end{lstlisting}

When a modal opens, push a fresh \code{FocusManager}. While the
stack is non-empty, route keys to the top of the stack. When the
modal closes, pop.

The visible UI also follows the stack: only the top modal's
elements are interactive; everything below is dimmed.

\begin{idea}
Modal focus has been a recurring source of bugs in
retained-mode GUIs (``why doesn't Tab work in this dialog?'').
Modelling it as a stack of focus managers is a couple of
\code{push\_back} / \code{pop\_back} calls; the runtime stays
unaware.
\end{idea}

\section{Pitfalls}

\begin{warning}
\textbf{Focus on a non-rendered element.} If \code{view} hides
the focused element (visibility, conditional), keystrokes go
nowhere. Always reconcile after the model changes: if the
focused id is no longer rendered, advance focus.
\end{warning}

\begin{warning}
\textbf{Stable IDs.} If you regenerate focus IDs every frame,
focus moves randomly. \code{register\_target} should be called
once, in \code{init} or when the focusable region is created.
\end{warning}

\begin{warning}
\textbf{Capturing focus during scroll.} If a focused element
scrolls out of view, the user can no longer see what they are
typing into. Either keep the focused element on-screen by
auto-scrolling, or visually indicate that focus is off-screen.
\end{warning}

\begin{warning}
\textbf{Forgetting to update tab order.} Adding a new focusable
region without inserting it into the order is a frequent
oversight. Centralise the tab order in one place rather than
sprinkling \code{register\_target} calls.
\end{warning}

\section*{Workshop}

\begin{enumerate}[leftmargin=*]
  \item Add a third focusable region to the two-counter app from
    Chapter~\ref{ch:elm}. Make Tab cycle through all three.
  \item Add a modal that opens on \code{?} and shows a help
    screen. While open, all keys should go to the modal.
  \item Read \file{maya/include/maya/core/focus.hpp}. Identify
    the data members of \code{FocusManager} and the operations
    that mutate them.
\end{enumerate}

\begin{recap}
Focus is state, not a framework property. A \code{FocusManager}
tracks IDs and order; \code{view} reads focus to choose styling;
\code{update} routes input by inspecting focus. Modals stack
focus managers. Mouse clicks update focus through a hit-test
mapping. No magic, no global pointer.
\end{recap}
